\documentclass[10pt,a4paper]{article}
\usepackage{authblk}
\renewcommand\Authfont{\footnotesize}
\renewcommand\Affilfont{\footnotesize\itshape}
\usepackage[margin=0.65in,top=0.6in,bottom=0.6in]{geometry}
\usepackage{amsmath,amssymb,amsfonts}
\usepackage{booktabs}
\usepackage{microtype}
\usepackage[hidelinks]{hyperref}

% Compact section formatting
\makeatletter
\renewcommand{\section}{\@startsection{section}{1}{\z@}{-1.0ex \@plus -0.4ex \@minus -.2ex}{0.4ex \@plus .2ex}{\normalsize\bfseries}}
\makeatother

\setlength{\parindent}{0pt}
\setlength{\parskip}{3pt}
\pagestyle{empty}

\begin{document}

% \begin{center}
% {\large\textbf{Inferring Latent Attacker Distributions from Time-Series Observations}}\\[0.25em]
% % {\footnotesize Research Update \quad\textbar\quad October 17, 2026 \quad\textbar\quad BATADAL Benchmark (Tank~1, $\sigma_{\text{train}} = 1.2085\,\text{m}$)}

% % {\footnotesize Aditya Hegde (Graduate Student), Saptarshi Sengupta (PI), San Jose State Unviersity}

% \end{center}



\begin{center}
{\large\textbf{Inferring Latent Attacker Distributions from Time-Series Observations}}\\[0.4em]
{\footnotesize NAIRR Pilot Start-Up Two-Month Update \quad\textbar\quad October 19, 2026 \quad\textbar\quad Request ID: \texttt{NAIRR260371}}\\[0.4em]
{\footnotesize
\begin{tabular}{c@{\hspace{3em}}c}
  {Aditya Hegde} (Graduate Student) & {Saptarshi Sengupta} (PI) \\
  \textit{Department of Computer Engineering}         & \textit{Department of Computer Science} \\
  \textit{San Jos\'e State University}     & \textit{San Jos\'e State University} \\
\end{tabular}}
\end{center}



\vspace{-0.5em}

\section{Problem Formulation \& Causal Filter Breakdown}
In industrial Cyber-Physical Systems (CPS), adversarial additive spoofing corrupts sensor telemetry at discrete timesteps $t \in \{1, \dots, T\}$ according to $y_t = x_t + \delta_t$, where $x_t \in \mathbb{R}$ is the unobserved physical plant state and $\delta_t \in \mathbb{R}$ is the unknown attack perturbation ($\delta_t = 0$ during clean operations). Recovering both trajectories from telemetry $y$ alone is non-identifiable without side information or structural priors: any decomposition $\tilde{\delta}_t = y_t - \tilde{x}_t$ yields identical observations. Standard causal estimators (Persistence, Seasonal Naive) fail under sustained spoofing because corrupted measurements directly contaminate their autoregressive states, driving severe divergence (Persistence median NRMSE $0.8616$; Seasonal Naive $1.4802$).

\section{Decoupled Inversion via Total Variation MAP}
To restore identifiability without ground truth at test time, we decouple the attacked channel from its corrupted history using uncompromised companion channels $z_t \in \mathbb{R}^d$ (downstream flows, pump statuses, adjacent pressures). A surrogate model yields clean state estimates $\hat{x}_{\text{surr}, t} = f_\theta(z_t)$ with residual variance $\sigma_\epsilon^2 = \mathbb{E}[(x_t - \hat{x}_{\text{surr}, t})^2]$. Exploiting piecewise-constant attack structure ($\sum_{t=1}^{T-1} |\delta_{t+1} - \delta_t| \ll T$), the perturbation trajectory is recovered via 1D Total Variation (TV) MAP optimization equivalently, the fused lasso signal approximator \cite{fused,rof}:
\begin{equation}
\min_{\delta \in \mathbb{R}^T} \;\; \frac{1}{2\sigma_\epsilon^2} \sum_{t=1}^T \left( y_t - \hat{x}_{\text{surr}, t} - \delta_t \right)^2 + \lambda_{\text{TV}} \sum_{t=1}^{T-1} |\delta_{t+1} - \delta_t| + \lambda_{\ell_1} \sum_{t=1}^T |\delta_t|
\end{equation}
The TV and $\ell_1$ penalties collapse high-frequency surrogate residual flutter into flat plateaus, recovering the plant state as $\hat{x}_t = y_t - \hat{\delta}_t$.

\section{Neural Surrogate Benchmarks}
Attacks are injected open-loop on the single Tank~1 channel of hourly telemetry and evaluated across 120 deterministic scenarios on the sealed BATADAL \cite{batadal} split (6h--48h durations, $0.5\sigma$--$3.0\sigma_{\text{train}}$ amplitudes), StepMAP with a linear surrogate achieves median NRMSE of {0.2031} (76.4\% error reduction over Persistence). Upgrading to a \texttt{DilatedTCNSurrogate} \cite{tcn} (48h context, trained on TACC Vista GH200; validation RMSE $0.0678\,\text{m}$, $R^2 = 0.9968$, compressing residual spread by 82.4\%) reduces median NRMSE to {0.1605} ($-21.0\%$ vs.\ Linear StepMAP; $-81.4\%$ vs.\ Persistence). For 48h attacks, error drops from {0.1689 to 0.0766} ($-54.6\%$), while weak $0.5\sigma$ attack NRMSE drops from {0.3832 to 0.2935} ($-23.4\%$) by lowering the surrogate noise floor.

\vspace{0.4em}
\begin{table}[htbp]
\centering
\small
\setlength{\tabcolsep}{4.5pt}
\renewcommand{\arraystretch}{1.05}
\begin{tabular}{lcccc}
\toprule
\textbf{Estimator} & \textbf{Median NRMSE} & \textbf{Mean NRMSE} & \textbf{Median Amp.\ Err.\ ($\sigma$)} & \textbf{Sign Acc.\ ($\ge 1\sigma$)} \\
\midrule
Persistence Baseline        & 0.8616 & 1.1672 & 0.6254 & 92.5\% \\
Seasonal Naive Baseline     & 1.4802 & 1.7063 & 1.2847 & 77.5\% \\
Companion Ridge Regression  & 0.3087 & 0.3650 & 0.0883 & 100.0\% \\
StepMAP (Linear Surrogate)  & 0.2031 & 0.2575 & 0.1441 & 100.0\% \\
StepMAP (Neural TCN) & {0.1605} & {0.2011} & {0.1446} & {100.0\%} \\
\bottomrule
\end{tabular}
\end{table}

\section{Next Steps}
In the remaining month of the allocation, we will explore multi-channel attacks and test robustness when companion channels are partially compromised. Consistent with NAIRR Pilot open-science requirements, we will release the implementation, trained surrogate and evaluation suite. 

\section{Acknowledgement}
The authors acknowledge the National Artificial Intelligence Research Resource (NAIRR) Pilot and the Texas Advanced Computing Center (TACC) Vista system for contributing to this research result.


\vspace{-0.5em}
{\footnotesize
\begin{thebibliography}{9}\setlength{\itemsep}{0pt}
\bibitem{fused} R.~Tibshirani, M.~Saunders, S.~Rosset, J.~Zhu, K.~Knight, ``Sparsity and Smoothness via the Fused Lasso,'' \emph{J.\ R.\ Stat.\ Soc.\ B}, 67(1):91--108, 2005.
\bibitem{rof} L.~I.~Rudin, S.~Osher, E.~Fatemi, ``Nonlinear Total Variation Based Noise Removal Algorithms,'' \emph{Physica D}, 60:259--268, 1992.
\bibitem{batadal} R.~Taormina \emph{et al.}, ``Battle of the Attack Detection Algorithms: Disclosing Cyber Attacks on Water Distribution Networks,'' \emph{J.\ Water Resour.\ Plann.\ Manage.}, 144(8):04018048, 2018.
\bibitem{tcn} S.~Bai, J.~Z.~Kolter, V.~Koltun, ``An Empirical Evaluation of Generic Convolutional and Recurrent Networks for Sequence Modeling,'' arXiv:1803.01271, 2018.
\end{thebibliography}}


\end{document}
